\documentclass[10pt,a4paper]{article}

% Shared presentation layer for the standalone R0--R7 development documents.
% Method equations remain authoritative in the canonical idea sketch.
\usepackage[a4paper,margin=18mm,headheight=14pt]{geometry}
\usepackage{fontspec}
\usepackage{xeCJK}
\xeCJKsetup{CJKspace=true}
\setmainfont{DejaVu Serif}
\setsansfont{DejaVu Sans}
\setmonofont{DejaVu Sans Mono}
\setCJKmainfont{Noto Serif CJK KR}
\setCJKsansfont{Noto Sans CJK KR}
\setCJKmonofont{Noto Sans Mono CJK KR}

\usepackage{amsmath,amssymb}
\usepackage{booktabs,longtable,tabularx,array}
\usepackage{enumitem}
\usepackage{xcolor}
\usepackage{hyperref}
\usepackage{fancyhdr}
\usepackage{microtype}
\usepackage[most]{tcolorbox}

\definecolor{CoreBlue}{HTML}{1F4E79}
\definecolor{CoreLight}{HTML}{EAF3F8}
\definecolor{StopRed}{HTML}{9C0006}
\definecolor{StopLight}{HTML}{FCE8E6}
\definecolor{GateGreen}{HTML}{38761D}
\definecolor{GateLight}{HTML}{EAF4E2}
\definecolor{DecisionOrange}{HTML}{B45F06}
\definecolor{DecisionLight}{HTML}{FFF2CC}

\hypersetup{
  colorlinks=true,
  linkcolor=CoreBlue,
  urlcolor=CoreBlue
}

\setlength{\parindent}{0pt}
\setlength{\parskip}{0.3em}
\setlength{\emergencystretch}{3em}
\setlist[itemize]{leftmargin=1.45em,itemsep=0.1em,topsep=0.2em}
\setlist[enumerate]{leftmargin=1.75em,itemsep=0.14em,topsep=0.22em}
\renewcommand{\arraystretch}{1.15}
\newcolumntype{L}[1]{>{\raggedright\arraybackslash}p{#1}}
\newcommand{\code}[1]{{\small\nolinkurl{#1}}}
\newcommand{\must}[1]{\textcolor{CoreBlue}{\textbf{#1}}}
\newcommand{\haltitem}[1]{\textcolor{StopRed}{\textbf{#1}}}
\newcommand{\passitem}[1]{\textcolor{GateGreen}{\textbf{#1}}}
\newcommand{\decisionitem}[1]{\textcolor{DecisionOrange}{\textbf{#1}}}

\newenvironment{contractbox}[1]
  {\begin{tcolorbox}[colback=CoreLight,colframe=CoreBlue,title={#1},fonttitle=\bfseries,before upper=\raggedright]}
  {\end{tcolorbox}}
\newenvironment{decisionbox}[1]
  {\begin{tcolorbox}[colback=DecisionLight,colframe=DecisionOrange,title={#1},fonttitle=\bfseries,before upper=\raggedright]}
  {\end{tcolorbox}}
\newenvironment{stopbox}[1]
  {\begin{tcolorbox}[colback=StopLight,colframe=StopRed,title={#1},fonttitle=\bfseries,before upper=\raggedright]}
  {\end{tcolorbox}}


\pagestyle{fancy}
\fancyhf{}
\lhead{Wind3DGS 개발 문서}
\rhead{R4 Held-out Global}
\cfoot{\thepage}

\title{R4 개발 계약\\Narrow-domain Held-out Global}
\author{Wind3DGS}
\date{2026-08-22}

\begin{document}
\maketitle

\begin{contractbox}{문서 역할}
이 문서는 object-disjoint narrow domain에서 learned Global response를 검증하고
Gate B 이후 사용할 단 하나의 Best Global checkpoint를 동결하기 위한 실행 계약이다.
Canonical idea sketch가 방법과 claim을 소유하며, 본 문서는 split, model selection,
held-out evaluation과 실패 routing을 소유한다.
Canonical \code{eq:rd-setup-input}, \code{eq:rd-response-package},
\code{eq:rd-global-response}, \code{eq:rd-wind-pullback}과
\code{eq:rd-exact-zoh}를 이름으로 참조한다.
\end{contractbox}

\tableofcontents

\section{목적과 소유권}

R4의 목적은 target mesh, target dynamic video와 per-target fine-tuning 없이 static 3DGS와
typed metric/material/attachment 입력만으로 held-out object의 passive Global response package를
예측할 수 있는지 검증하는 것이다. R4는 다음을 소유한다.

\begin{itemize}
  \item source-object grouped train/dev/test split와 test-open 절차
  \item R3가 hash한 route/config 후보 안의 Global rank/capacity와 checkpoint dev selection
  \item R2b가 동결한 response loss와 rollout schedule의 read-only 소비
  \item direct next-deformation, direct-set response 및 fixed retrieval baseline과의 공정 비교
  \item material/typed-control ordering, long rollout와 independent-resplat worst-case 평가
  \item Gate B Boolean 판정과 R5에 넘길 \code{BestGlobalCheckpoint} freeze
\end{itemize}

R4는 Local residual, conditional selector 또는 appearance branch를 학습하지 않는다.
Gate B가 닫히기 전에는 R5--R7 결과를 보고 Global config를 다시 선택하지 않는다.

\section{진입 입력}

\begin{longtable}{L{0.28\linewidth}L{0.60\linewidth}}
\toprule
입력 & 요구 내용\\
\midrule
\endfirsthead
\toprule
입력 & 요구 내용\\
\midrule
\endhead
\code{DomainContract} & supported strip/flag family, metric/material/attachment/wind/control envelope와 OOD 규칙\\
\code{ObjectGroupedSplit} & R0가 동결한 object identity와 mesh/GS/material/wind의 leakage-free grouping key\\
\code{GateAReport} & Gate A\(_T\) pass와 A\(_M\)/A\(_L\) route 판정, threshold/config hash\\
\code{R3RouteConfig} & latent 또는 direct-set main architecture, fixed-family 제한과 invalidation hash\\
\code{R3RouteCheckpoint} & 선택 route의 dev checkpoint와 R4 initialization/training provenance\\
\code{TeacherPackage} & converged teacher package collection, canonical probe maps, external work와 sequence metadata\\
\code{TeacherSpectrumStatus} & R1에서 동결한 extraction config/hash와 stable-peak available/unavailable status\\
\code{R2bHandoff} & R2b가 채택한 response loss, rollout schedule/fallback branch, horizon과 compute-matching hash\\
\code{MultiResplatGroupManifest} & R3의 independent GS variants, eligibility와 density-claim 상태\\
\code{GlobalConfigSet} & R3가 hash한 route architecture와 유한 Global rank/pole/capacity/compute 후보; R2b schedule은 포함하지 않음\\
\code{MetricRegistry} & response, material ordering, control, stability와 system metric 정의\\
\code{ThresholdRegistry} & Gate B에 사용할 dev-frozen threshold와 Boolean 식; 미동결이면 R4 종료 불가\\
\bottomrule
\end{longtable}

R0--R3의 unit, common-probe, evaluator, measure와 eligibility hash가 바뀌면 기존 R4 run을
자동 승계하지 않는다. Compatibility manifest가 같지 않으면 새 실험 lineage를 만든다.

\section{이미 동결된 계약}

\subsection{Object split와 inference boundary}

\begin{itemize}
  \item 같은 source object의 mesh, 모든 independent GS variant, material과 wind sequence는
  하나의 train/dev/test group에 둔다.
  \item Test object에는 원칙적으로 2--3개의 declared independent GS density/seed variant를 둔다.
  R3에서 density claim을 철회했다면 fixed-family variant만 사용하고 그 범위를 명시한다.
  \item Test target의 mesh, dynamic RGB/video, teacher trajectory와 correspondence는 setup input,
  fine-tuning 또는 model selection에 사용하지 않는다.
  \item Test-open 전 split, checkpoint selection rule, metric, threshold, control range와 success denominator를 hash한다.
  \item Object 수가 작으면 exploratory narrow-domain evidence로 보고 population-level 일반화 claim을 하지 않는다.
\end{itemize}

\subsection{Gate A route ownership}

\begin{itemize}
  \item Gate A\(_T\)가 pass가 아니면 R4를 시작하지 않는다.
  \item Gate A\(_L\) pass이면 latent-relation model이 main route이고 capacity-matched direct-set은 baseline이다.
  \item Gate A\(_L\) fail이면 direct-set response가 main route이며 latent model은 diagnostic/baseline으로만 남긴다.
  \item Gate A\(_M\) fail이면 어느 encoder route든 R3에서 동결한 fixed reconstruction family에만 적용하고
  density-invariance claim을 사용하지 않는다.
  \item R4 dev/test 결과를 보고 A\(_M\)/A\(_L\) route를 다시 선택하지 않는다.
\end{itemize}

\subsection{Global package와 evaluator}

Global model은 static GS로부터 cached response package를 한 번 출력한다. Frame evaluator는
frame-start traction sample과 exact-ZOH oscillator update를 사용하며 network가 매 frame next pose를
다시 예측하지 않는다. Attachment row, mass Gram, positive angular frequency/modal damping ratio와
package hash validation은 construction/fixture로 통과해야 한다.

Global은 R4 동안 always-on이며 Local field/state가 없다. Base field/pole은 immutable하다.
Typed controls는 package 허용 범위에서만
\[
  W_t^{\mathrm{eff}}=g_W(t)W_t,\qquad
  \Omega_G^{\mathrm{eff}}=\sqrt{g_K}\,\Omega_G^{\mathrm{base}},\qquad
  Z_G^{\mathrm{eff}}=g_D Z_G^{\mathrm{base}}
\]
로 파생한다. \(g_K,g_D\) 변경은 rest reset과 새 derived evaluator hash를 요구하며,
material/attachment 변경은 setup rerun과 새 package를 요구한다.

\subsection{Model selection privilege}

R4의 선택 권한은 hash된 \code{R3RouteConfig}/\code{GlobalConfigSet} 안의 rank, capacity와 checkpoint에만 있다.
R2b의 loss weight, curriculum/fallback, horizon과 compute-matching config는 read-only로 소비한다.
Best Global 선택 후에는 다음을 하나의 immutable identity로 고정한다.

\begin{itemize}
  \item network architecture와 parameter checkpoint
  \item Global rank/pole 수, token/patch config와 package schema
  \item loss, optimizer, schedule, stopping epoch와 random-seed policy
  \item common-probe, evaluator, metric/threshold registry와 supported control interval
  \item setup/runtime precision, batch와 deployment profiling config
\end{itemize}

Encoder family, capacity-matching 조건 또는 \code{GlobalConfigSet} 밖 후보가 필요하면 R4에서 search를 넓히지 않고
Gate A\(_L\) evidence를 무효화해 R3부터 새 lineage로 실행한다. R2b schedule을 바꿔야 하면 R2를 무효화한다.
Test 결과나 R5 Local 결과를 보고 Best Global을 다시 선택하거나 per-asset pole/force scale을 보정하지 않는다.

\section{구현 체크리스트}

\subsection{Split와 dataset audit}

\begin{itemize}
  \item[$\square$] Stable source-object fingerprint와 group ID 생성
  \item[$\square$] Mesh/GS/material/wind 파일에서 train/dev/test 교차 중복 검사
  \item[$\square$] Aspect ratio, silhouette, material preset, metric scale와 attachment-region 분포표 생성
  \item[$\square$] Test에는 unseen phase/two-tone 또는 reversed traveling-gust composition 포함
  \item[$\square$] 학습보다 긴 zero-wind recovery/free-decay sequence를 dev/test에 고정
  \item[$\square$] Split manifest와 test seal hash를 학습 전에 저장
\end{itemize}

\subsection{Global training과 model selection}

\begin{itemize}
  \item[$\square$] \code{R3RouteConfig}의 main encoder route와 fixed-family 제한을 그대로 load하고 compatibility hash 검사
  \item[$\square$] \code{GlobalConfigSet} 안의 predeclared Global rank/capacity/config만 동일 train/dev compute 규칙으로 실행
  \item[$\square$] Encoder family와 latent/direct-set capacity-matching identity가 R3 hash와 같은지 검사
  \item[$\square$] Direct next-deformation과 fixed mean/nearest-response baseline 실행
  \item[$\square$] A\(_L\) pass에서는 capacity-matched direct-set을 baseline으로, fail에서는 main route로 실행
  \item[$\square$] Common-probe response loss의 displacement/velocity/FRF/rollout normalization과 reduction 확인
  \item[$\square$] \code{TeacherSpectrumStatus}와 \code{R2bHandoff}를 대조해 R2b가 동결한
  period 또는 fixed-seconds branch를 그대로 사용
  \item[$\square$] R4 결과로 curriculum/fallback와 final horizon을 다시 선택하지 않음
  \item[$\square$] Spatial/temporal nonconvergence나 final horizon divergence를 schedule fallback으로 숨기지 않음
  \item[$\square$] Dev-only selection report가 단 하나의 winner와 deterministic tie-break를 반환
  \item[$\square$] Winner checkpoint/package/config와 모든 dependency hash를 freeze
  \item[$\square$] Best Global report에 latent/direct-set main route와 Gate A claim restriction을 명시
\end{itemize}

\subsection{Held-out evaluation}

\begin{itemize}
  \item[$\square$] Test seal 확인 후 frozen Best Global을 한 번만 load하여 모든 object/variant 실행
  \item[$\square$] No target mesh/video/fine-tuning 로그와 package setup input을 감사
  \item[$\square$] Impulse, step/recovery, multi-sine/chirp와 traveling gust의 velocity/phase/PSD 평가
  \item[$\square$] Soft/base/stiff material amplitude/frequency ordering 측정
  \item[$\square$] \(g_W,g_K,g_D\)를 dev-frozen interval에서 별도 rest-reset rollout으로 sweep
  \item[$\square$] Expected monotonicity, phase continuity, aero-off free-decay와 damping ordering 검사
  \item[$\square$] Final supported physical-time horizon의 NaN/divergence와 attachment residual 검사
  \item[$\square$] Independent-resplat worst case, method success rate와 excluded denominator 보고
  \item[$\square$] Object-outer/package-within-object nested paired statistics와 low-\(N\) limitation 보고
\end{itemize}

\subsection{Best Global hand-off}

\begin{itemize}
  \item[$\square$] \code{BestGlobalCheckpoint}를 read-only artifact로 export
  \item[$\square$] R5가 사용할 Global output과 canonical-probe teacher residual 생성 API version 고정
  \item[$\square$] Rank/capacity/config/checkpoint 재선택 금지를 compatibility validator로 검사
  \item[$\square$] R5용 larger-Global sweep 후보와 matched-p95 profiling protocol은 별도 baseline registry로 고정
\end{itemize}

\section{Open Design Decisions}

\begin{decisionbox}{Gate B는 아직 수치적으로 미완성}
Gate B의 자연어 목표는 정해졌지만 primary metric, numeric threshold, aggregation과 Boolean AND/OR 식은
아직 구현 계약으로 확정되지 않았다. \code{ThresholdRegistry}를 dev 결과로 동결하기 전에는
held-out 결과가 좋아 보여도 Gate B pass를 선언할 수 없다.
\end{decisionbox}

\begin{enumerate}
  \item \decisionitem{Object split 규모:} train/dev/test object 수와 family별 최소 coverage를 확정한다.
  \item \decisionitem{Best Global selection:} hash된 \code{GlobalConfigSet} 안에서 사용할 primary scalar,
  Pareto rule, stability veto와 deterministic tie-break를 정한다.
  \item \decisionitem{Gate B Boolean registry:} baseline 대비 허용/개선 threshold, material/control ordering,
  divergence veto, object success ratio와 worst-variant 조건을 정확한 Boolean 식으로 정한다.
  \item \decisionitem{통계 집계:} object-outer mean/median/worst, package-within-object effect,
  confidence interval과 multiple-seed reduction 방식을 정한다.
  \item \decisionitem{Control range:} \(g_W,g_K,g_D\)의 dev-frozen interval과 monotonicity allowance를 정한다.
\end{enumerate}

\section{산출물}

\begin{longtable}{L{0.30\linewidth}L{0.58\linewidth}}
\toprule
산출물 & 필수 내용\\
\midrule
\endfirsthead
\toprule
산출물 & 필수 내용\\
\midrule
\endhead
\code{DataSplitManifest} & source groups, train/dev/test, leakage audit, test seal와 split hash\\
\code{GlobalSearchRegistry} & 유한 config/seed/compute 후보와 실행 상태\\
\code{ModelSelectionReport} & dev primary/secondary metric, veto, tie-break와 winner 근거\\
\code{BestGlobalCheckpoint} & immutable weights, latent/direct-set route, config, schema, dependency와 compatibility hash\\
\code{GlobalResponsePackageSet} & Best Global model로 declared asset마다 compile한 immutable \code{ResponsePackage}, model/config/evaluator dependency hash\\
\code{HeldoutGlobalReport} & object/variant별 response, control, stability, success denominator와 baselines\\
\code{GateBReport} & frozen threshold registry 입력, Boolean subcondition과 최종 pass/fail\\
\code{R5GlobalHandoff} & residual extraction API, teacher/probe identity와 larger-Global baseline registry\\
\bottomrule
\end{longtable}

\section{필수 fixture와 metric}

\subsection{필수 fixture}

\begin{enumerate}
  \item \textbf{Split leakage:} source fingerprint와 모든 derived asset의 group consistency
  \item \textbf{Package validity:} SI unit, rank, positive pole, mass Gram, attachment zero와 serialization replay
  \item \textbf{Held-out response:} impulse, step/recovery, frequency probe와 traveling gust
  \item \textbf{Material ordering:} soft/base/stiff의 amplitude, dominant frequency와 decay
  \item \textbf{Typed control:} wind/stiffness/damping gain의 monotonicity, phase continuity와 rest-reset identity
  \item \textbf{Long horizon:} final supported time, zero-wind recovery와 aero-off free decay
  \item \textbf{Multi-resplat worst case:} object 안의 모든 declared eligible variant
  \item \textbf{System profile:} setup, package memory, force/update/transport/render p50/p95/p99
\end{enumerate}

\subsection{Metric 층위}

\begin{itemize}
  \item \textbf{Primary 후보:} normalized common-probe velocity, phase와 target-band response error
  \item \textbf{필수 veto:} final-horizon divergence, attachment violation, test privilege 위반
  \item \textbf{Ordering:} material 및 typed control의 expected amplitude/frequency/decay ordering
  \item \textbf{Robustness:} worst independent-resplat degradation와 object-level success rate
  \item \textbf{Secondary:} displacement/tip, PSD, external work, setup/runtime memory와 latency
\end{itemize}

Primary 확정과 Gate B 수치식은 Open Design Decision이며 test-open 전에 registry로 동결한다.

\section{종료 조건과 실패 경로}

\begin{contractbox}{R4 종료 조건}
Leakage-free split, dev-only model selection, immutable Best Global export와 sealed-test report가 완성되고,
frozen \code{ThresholdRegistry}가 Gate B Boolean 결과를 재현하면 R4를 종료한다.
Gate B pass일 때만 해당 Best Global checkpoint를 R5의 frozen residual owner로 넘긴다.
\end{contractbox}

\begin{itemize}
  \item \passitem{Gate B pass:} checkpoint/config/rank를 고정하고 R5 residual atlas로 진행한다.
  \item \haltitem{Supported-family partial failure:} dev 규칙에 따라 family를 더 좁히거나
  per-family compiler로 claim을 축소한다. Test 결과를 보고 같은 split에서 재선택하지 않는다.
  \item \haltitem{Held-out response failure:} target-mesh-free amortized learned-response claim을 중단한다.
  Dataset 확대나 Local 추가로 Global 실패를 덮지 않는다.
  \item \haltitem{Divergence/decay-order failure:} final supported horizon에서 한 건이라도 발생하면
  Gate B를 실패로 처리하고 evaluator/package/domain owner를 진단한다.
  \item \haltitem{No density claim:} R3 A\(_M\) 실패 상태라면 fixed reconstruction family 결과만 주장하며
  R4의 좋은 결과로 density invariance를 복원하지 않는다.
  \item \haltitem{No latent claim:} R3 A\(_L\) 실패 상태라면 direct-set을 main route로 유지하고
  R4의 좋은 결과로 latent-anchor/relation contribution을 복원하지 않는다.
  \item \haltitem{Threshold 미동결:} qualitative success로 표시할 수는 있으나 Gate B와 Best Global hand-off는 보류한다.
\end{itemize}

\end{document}
